\chapter{Hasard et probabilités}
\label{exercices:11}
\section*{Exercice 11.1 -- Une urne}
Une urne contient $3$ boules rouges, $2$ bleues et $5$ vertes, indiscernables au toucher. Calculer la probabilité de chaque couleur, puis de « ne pas obtenir vert ». Les couleurs sont-elles équiprobables ?

\section*{Exercice 11.2 -- Deux pièces}
Deux pièces équilibrées sont lancées indépendamment. Énumérer les issues ordonnées. Calculer les probabilités d'obtenir exactement une fois pile, au moins une fois pile et deux fois le même côté.

\section*{Exercice 11.3 -- Avec puis sans remise}
Une urne contient les boules $R_1,R_2,B$. On en tire deux successivement. Calculer la probabilité d'obtenir deux rouges avec remise, puis sans remise. Construire un tableau ou un arbre pour justifier.

\section*{Exercice 11.4 -- Deux dés et simulation}
Deux dés équilibrés indépendants sont lancés. Calculer la probabilité d'une somme égale à $9$. Une simulation obtient $114$ succès sur $1000$ essais : donner sa fréquence et commenter. Pourquoi simuler directement un entier entre $2$ et $12$ serait-il incorrect ?
